\uuid{oOEP}
\titre{Équation}
\auteur{bodin ; hanani ; mzari}
\organisation{exo7}
\licence{CC BY-NC-SA 4.0}
\datecreate{2020-01-20}
\niveau{L1}
\module{Analyse}
\chapitre{Fonctions circulaires et hyperboliques inverses}
\sousChapitre{Autre}
\difficulte{3}
\theme{Fonctions usuelles, équation}
\exo7qcmid{lille-1.fonctions-usuelles.022}

\contenu{
\begin{question}
\qclassification{126.00}
\qtags{motcle=équation}
Soit $(S)$ le système d'équations : $\left\{\begin{array}{ccl}2^x&=&y^2\\2^{x+1}&=&y^{2+x} \end{array}\right.$.  On note $E$ l'ensemble des $(x,y)$ qui vérifient $(S)$. Quelles sont les assertions vraies ?

\begin{answers}
    \bad {$(S)$ est défini  pour tout $ (x,y) \in \Rr ^2$.}

    \bad{Le cardinal de $E$ est $1$.}

    \good{Le cardinal de $E$ est $2$.}

    \bad{Le cardinal de $E$ est $4$.}
\end{answers}
\begin{explanations}
$(S)$ est défini pour tout $x\in \Rr$ et $y>0$.

Soit $x\in \Rr$ et $y>0$, $(S) \Leftrightarrow \left\{\begin{array}{ccl}y&=&2^{\frac{x}{2}} \\2^{x+1}&=&2^{x+\frac{x^2}{2}} \end{array}\right.  \Leftrightarrow \left\{\begin{array}{ccl}y&=&2^{\frac{x}{2}} \\x^2&=&2 \end{array}\right. $.

Donc $E = \{ (\sqrt 2, \sqrt 2^{\sqrt 2})\,  ; (-\sqrt 2, \sqrt 2^{-\sqrt 2})\}$ et donc le cardinal de $E$ est $2$.
\end{explanations}
\end{question}
}
